\documentclass[11pt]{article}
\usepackage[hmargin=2.5cm,vmargin=2.5cm]{geometry}
\usepackage{amsmath,amssymb}

\title{The technology growth rate $1+\gamma^{a}$ vs.\ the output growth rate $g$}
\author{Spending-efficiency model --- internal note}
\date{\today}

\begin{document}
\maketitle

\noindent
\textbf{Finding.} The model's technology block (adoption law of motion, the value
functions $\mathcal V$ and $\mathcal J$, and the $S$ first-order condition) is detrended
by \texttt{(1+gammaa)}, \emph{not} by the output growth rate $g$. The paper's appendix
originally wrote those same equations with $g$, which does not match the solved model.
This note records why the model is right and what was fixed.

\section{$\gamma^{a}$ is derived from $g$}

\texttt{gammaa} is not a free parameter. The code sets
\begin{equation}
    \texttt{gammaa = g\^{}((1-alpha)/(vartheta-1)) - 1},
    \qquad\text{i.e.}\qquad
    1+\gamma^{a} = g^{\,(1-\alpha)/(\vartheta-1)} .
\end{equation}
With the calibration $\alpha = 0.3$ and $\vartheta = 1.35$ the exponent is exactly
$0.7/0.35 = 2$, so
\begin{equation}
    1+\gamma^{a} = g^{2}.
\end{equation}
Technology therefore grows at the \emph{square} of the output growth factor: $1.008$
versus $1.004$ (AE) and $1.015$ versus $1.0075$ (EMDE).

\section{Why technology must grow faster than output}

Technology enters production only through the love-of-variety term $A_t^{\vartheta-1}$:
\begin{equation}
    Y_t = A_{t-1}^{\vartheta-1} K_{t-1}^{\alpha} N_t^{1-\alpha}
          \left(K^{GI}_{t-1}\right)^{\alpha_G} - \Theta .
\end{equation}
On the balanced-growth path output and capital grow at $g$ and labor is stationary, so
\begin{equation}
    g^{1-\alpha} = \left(1+\gamma^{a}\right)^{\vartheta-1}
    \qquad\Longrightarrow\qquad
    1+\gamma^{a} = g^{\,(1-\alpha)/(\vartheta-1)},
\end{equation}
exactly the code's formula (it ignores the small $\alpha_G$ term). Because the technology
exponent $\vartheta-1 = 0.35$ is below the labor share $1-\alpha = 0.7$, the technology
stock must grow faster than output for TFP to lift output at rate $g$. Equivalently, the
real trend is $\Lambda_t = A_t^{(\vartheta-1)/(1-\alpha)} = A_t^{1/2}$: real quantities
grow at $g$, the technology stock at $g^{2}$.

\section{Consequence and fix}

A single trend at $g$ for everything is \emph{not} innocuous: it would make the
production function fail to balance (with $A$ growing at $g$, output would grow at
$g^{\vartheta-1+\alpha} \neq g$). The model correctly runs two trends---$g$ for the real
block, $1+\gamma^{a}$ for technology---so \texttt{(1+gammaa)} in the adoption block is the
right factor, and the paper's $g$ there was the error.

Fixed in the paper:
\begin{itemize}
    \item Appendix technology block---the adoption law of motion and the value/FOC
          equations now carry $1+\gamma^{a}$ in place of $g$.
    \item A sentence defines $1+\gamma^{a}$ and states
          $1+\gamma^{a} = g^{(1-\alpha)/(\vartheta-1)}$; Section~2.5 flags the two trends.
    \item The glossary gives \texttt{gammaa} the symbol $\gamma^{a}$, and the endogenous
          table's equations for $A$, $\mathcal V$, $\mathcal J$, $S$ use $1+\gamma^{a}$.
\end{itemize}
The model was not changed; this was a paper-side correction to match the code.

\end{document}
